% !TEX program = lualatex
\documentclass[a4paper,11pt]{ltjsarticle}

\usepackage{amsmath,amssymb}
\usepackage{geometry}
\usepackage{listings}
\usepackage{xcolor}
\geometry{margin=25mm}

\lstset{
  basicstyle=\ttfamily\small,
  columns=fullflexible,
  frame=single,
  keepspaces=true
}

\title{線形写像の核・像と次元定理\\
\large GAPを使った計算例}
\author{ecsm}
\date{}

\begin{document}

\maketitle

\section*{問題}

線形写像
\[
 f:\mathbb{R}^3\longrightarrow\mathbb{R}^2,\qquad
 f(x)=Ax,
\]
を
\[
 A=
 \begin{pmatrix}
  1&2&1\\
  2&1&3
 \end{pmatrix}
\]
で定める。

ここでは通常の線形代数の教科書と同様に、ベクトルを列ベクトルとして表す。

\begin{enumerate}
  \item $\ker f$ を求めよ。
  \item $\operatorname{Im}f$ を求めよ。
  \item
  \[
    \dim\mathbb{R}^3
    =
    \dim\ker f+\dim\operatorname{Im}f
  \]
  を確認せよ。
\end{enumerate}

\section*{1. 核を求める}

$Ax=0$ を解く。

\[
\begin{pmatrix}
1&2&1\\
2&1&3
\end{pmatrix}
\begin{pmatrix}
x\\y\\z
\end{pmatrix}
=
\begin{pmatrix}
0\\0
\end{pmatrix}.
\]

これは
\[
\begin{cases}
x+2y+z=0,\\
2x+y+3z=0
\end{cases}
\]
と同値である。

第2式から第1式の2倍を引くと
\[
-3y+z=0,
\]
したがって $z=3y$ である。また
\[
x+2y+3y=0
\]
より $x=-5y$ となる。

よって
\[
(x,y,z)=y(-5,1,3)
\]
であり、
\[
\boxed{\ker f=\operatorname{span}\{(-5,1,3)\}}.
\]

したがって
\[
\dim\ker f=1.
\]

\section*{2. 像を求める}

行列 $A$ の列ベクトルは
\[
\begin{pmatrix}1\\2\end{pmatrix},
\quad
\begin{pmatrix}2\\1\end{pmatrix},
\quad
\begin{pmatrix}1\\3\end{pmatrix}
\]
である。

第1列と第2列からなる行列
\[
\begin{pmatrix}
1&2\\
2&1
\end{pmatrix}
\]
の行列式は
\[
1\cdot1-2\cdot2=-3\neq0.
\]
したがって第1列と第2列は一次独立であり、
\[
\operatorname{Im}f=\mathbb{R}^2.
\]

よって
\[
\boxed{\dim\operatorname{Im}f=2}.
\]

\section*{3. 次元定理の確認}

以上から
\[
\dim\ker f+\dim\operatorname{Im}f
=1+2=3.
\]
一方、
\[
\dim\mathbb{R}^3=3.
\]
したがって
\[
\boxed{
\dim\mathbb{R}^3
=
\dim\ker f+\dim\operatorname{Im}f
}.
\]

これは線形写像に対する
\[
\boxed{\dim V=\dim\ker f+\dim\operatorname{Im}f}
\]
という次元定理の具体例になっている。

\section*{4. GAPで計算する}

ここでは、通常の線形代数とGAPでのベクトルの表現の違いに注意する。

数学ではベクトルを列ベクトルとして
\[
x=
\begin{pmatrix}
x_1\\x_2\\x_3
\end{pmatrix}
\]
と表し、線形写像は
\[
Ax
\]
と書く。

一方、GAPではベクトルを行リストとして扱う。
したがって、GAPでは
\[
[x_1,x_2,x_3]
\]
という形でベクトルを表し、対応する線形写像は転置行列を右から掛ける形
\[
[x_1,x_2,x_3]A^{\mathsf T}
\]
で表すことができる。

ここで
\[
A^{\mathsf T}
=
\begin{pmatrix}
1&2\\
2&1\\
1&3
\end{pmatrix}
\]
なので、
\[
[x_1,x_2,x_3]A^{\mathsf T}
=
[x_1+2x_2+x_3,\,
  2x_1+x_2+3x_3].
\]

これは数学での
\[
A
\begin{pmatrix}
x_1\\x_2\\x_3
\end{pmatrix}
\]
と同じ線形写像を表している。

\subsection*{4.1 行列をGAPに入力する}

まず、数学で使っている行列 $A$ をGAPに入力する。

\begin{lstlisting}
gap> A := [[1,2,1],[2,1,3]];;
\end{lstlisting}

GAPで行列のランクを計算すると、

\begin{lstlisting}
gap> RankMat(A);
2
\end{lstlisting}

となる。

したがって
\[
\operatorname{rank}(A)=2
\]
であり、
\[
\dim\ker f
=3-\operatorname{rank}(A)
=1
\]
である。

\subsection*{4.2 核の計算}

GAPでは、行列 $A$ の右核（right nullspace）を考えることで、
\[
xA^{\mathsf T}=0
\]
を満たす行ベクトル $x$ を求めることができる。

\begin{lstlisting}
gap> AT := TransposedMat(A);;
gap> NullspaceMat(AT);
[ [ -5, 1, 3 ] ]
\end{lstlisting}

したがって
\[
[-5,1,3]
\]
が核の基底に対応する。

これは数学での列ベクトル
\[
\begin{pmatrix}
-5\\1\\3
\end{pmatrix}
\]
に対応しているので、
\[
\boxed{\ker f
=\operatorname{span}\{(-5,1,3)\}}
\]
がGAPでも確認できる。

\subsection*{4.3 像の計算}

数学では $\operatorname{Im}f$ は $A$ の列空間である。

GAPでは、$A$ の列を行ベクトルとして扱うために転置を利用する。
すなわち、$A^{\mathsf T}$ の行空間を考える。

\begin{lstlisting}
gap> BaseMat(AT);
[ [ 1, 2 ], [ 2, 1 ] ]
\end{lstlisting}

この2本は
\[
(1,2),\qquad(2,1)
\]
に対応し、数学での像の基底
\[
\begin{pmatrix}1\\2\end{pmatrix},
\qquad
\begin{pmatrix}2\\1\end{pmatrix}
\]
を与える。

したがって
\[
\boxed{\operatorname{Im}f=\mathbb R^2}
\]
であり、
\[
\dim\operatorname{Im}f=2
\]
である。

\section*{5. 手計算とGAPの対応}

この例では、数学での列ベクトルとGAPでの行ベクトルを次のように対応させることができる。

\[
\begin{array}{c|c}
\text{線形代数} & \text{GAP}\\ \hline
x\in\mathbb R^3\text{（列ベクトル）}
&
[x_1,x_2,x_3]\text{（行リスト）}
\\[2mm]
Ax
&
xA^{\mathsf T}
\\[2mm]
\ker f
&
\texttt{NullspaceMat(TransposedMat(A))}
\\[2mm]
\operatorname{Im}f
&
A^{\mathsf T}\text{ の行空間}
\\[2mm]
\dim\operatorname{Im}f
&
\texttt{RankMat(A)}
\end{array}
\]

重要なのは、数学での定義そのものをGAPの表記に合わせて変更する必要はないということである。

数学では通常通り列ベクトルを使い、
\[
f(x)=Ax
\]
と考える。
GAPではその同じ線形写像を、行ベクトルと転置行列を使って
\[
xA^{\mathsf T}
\]
と表現すればよい。

\section*{6. 次元定理をGAPで確認する}

GAPで得られた結果は
\[
\operatorname{rank}(A)=2,
\qquad
\dim\ker f=1.
\]

また、$A$ は $2\times3$ 行列なので、
\[
\dim\operatorname{Im}f
=\operatorname{rank}(A)
=2.
\]

したがって
\[
\dim\ker f+\dim\operatorname{Im}f
=1+2=3
=\dim\mathbb R^3.
\]

よって、この具体例において次元定理
\[
\boxed{
\dim V
=
\dim\ker f+\dim\operatorname{Im}f
}
\]
が計算によって確認できた。

\section*{まとめ}

この例では、線形写像の核と像を手計算で求め、その結果をGAPでも確認した。

GAPを使う際には、数学で通常用いる列ベクトルと、GAPで扱われる行ベクトルの違いに注意する必要がある。
しかし、転置行列を用いれば両者を自然に対応させることができる。

GAPによる計算は証明そのものを置き換えるものではない。
一方で、具体例を計算することで、核・像・ランク・次元定理がどのように結び付いているかを確認することができる。

\end{document}
